\chapter{Architecture}

CharlotteOS is built on three kernel primitives: Capabilities, Endpoints, and
Memory Objects. Every other facility (completion queues, reply tokens, device
drivers) composes from these.

\section{The three primitives}

\subsection{Capabilities}

A capability is authority --- a per-address-space handle that names a kernel
object and carries rights. An application cannot fabricate a capability; it can
only receive one through the bootstrap contract or through IPC. Examples:
endpoint capabilities, connection capabilities, memory-object capabilities,
interrupt capabilities, MMIO-region capabilities.

\subsection{Endpoints}

An endpoint is a bounded server-side IPC receive object: it has an owning
domain, a queue capacity, and a lifecycle. Clients communicate through
\textbf{connection capabilities} minted from the endpoint. The kernel validates
and enqueues messages; the server receives them. Endpoints know nothing about
packet buffers, futures, or ownership --- they carry requests and replies.

\subsection{Memory Objects}

A memory object is page-backed storage with explicit ownership semantics. It can
be mapped into a domain's address space, moved between domains (ownership
transfer), lent for temporary read or write access, or copied. The MMU enforces
the access rights; Rust ownership complements but does not replace kernel
enforcement.

\section{Endpoint IPC}

\subsection{Scalar messages}

A scalar message carries an opcode, one 64-bit argument, and optionally a reply
token. It is suitable for control-plane operations: opcodes, handles, status
values.

\subsection{Memory-bearing messages}

A message can carry one or more memory-object attachments with explicit transfer
modes:

\begin{itemize}
    \item \textbf{Copy} --- the receiver gets a byte copy; the sender retains
        the original.
    \item \textbf{Move} --- ownership transfers to the receiver; the sender
        loses access.
    \item \textbf{BorrowRead} --- the receiver gets a temporary read-only
        mapping, revoked on reply or cancellation.
    \item \textbf{BorrowWrite} --- the receiver gets a temporary exclusive
        writable mapping, revoked on reply or cancellation.
\end{itemize}

\subsection{Connection delegation}

A connection capability can be delegated to another domain with attenuated
rights. This is the mechanism by which the userspace name service returns
service-level connections: it mints a connection with \texttt{SEND | CALL}
rights from a re-delegable connection it received at registration time. The
kernel mediates; the name service never sees raw endpoint identities.

\subsection{Deferred replies}

A server may retain a reply token and complete it later --- while continuing to
serve other requests on the same shard. This is essential for device-driven
operations: an interrupt completes a deferred read, the driver replies with the
received data, and no thread is dedicated to the wait.

\section{Service supervision}

A \textbf{supervisor} (kernel-side) creates protection domains, loads ELF
images, and delivers exactly one bootstrap capability per domain through a
config-page contract. The bootstrap connection is either an endpoint (for the
name service) or a connection to the name service (for ordinary services and
clients). For driver domains, additional device capabilities (MMIO region,
interrupt) are delivered through the same config page.

A userspace \textbf{name service} maps human-readable names to
re-delegable connection capabilities with instance generations. Re-registration
bumps the generation and closes the previous instance's connection; stale client
connections fail deterministically with \texttt{EndpointClosed}. The name
service supports access-key gating: a registration can specify a non-zero
access key, and lookups without the matching key are denied.

\section{Device capabilities}

Two first-class object types extend the capability model for userspace drivers:

\begin{itemize}
    \item \textbf{MmioRegion} --- a page-granular device register window that a
        driver maps into its own address space as Device-nGnRnE memory
        (user-accessible, execute-never, strongly ordered).
    \item \textbf{Interrupt} --- a routed interrupt source whose readiness is
        delivered to the driver's completion queue as a coalesced wake.
\end{itemize}

Grants are minted kernel-side by the supervisor and never through a syscall. A
driver receives exactly the MMIO window and interrupt it needs, with no access
to arbitrary physical memory or interrupt vectors.

Interrupt delivery follows a lock-free path: the IRQ dispatcher reads the
per-INTID route table atomically, masks the source via MMIO, atomically bumps
coalescing counters, and pushes a wake onto a deferred queue. The actual
\texttt{completion::wake} (which takes locks) is performed from thread context
by \texttt{yield\_lp} and the LP idle loop.

\section{Completion queues}

A completion queue (CQ) is a shared-memory ring buffer between the kernel
(producer) and userspace (consumer). Each entry is 32 bytes (\S8.2):
\texttt{operation: u64, cookie: u64, status: u32, flags: u32, result: i64}.
Entries are published with a release fence and a single head update; the
consumer advances the tail after reading. A non-lossy kernel backlog retains
entries when the ring is full --- terminal completions are never dropped.

Each address space owns a set of queues keyed by \texttt{CqId}. Queue 0 is the
default. Per-shard queues (one per sitas shard) enable targeted wakes, so a
cross-shard message wakes only the target shard's executor.

\texttt{CQ\_WAIT} blocks the caller until a completion entry is posted, an
explicit \texttt{CQ\_WAKE} arrives, or a deadline elapses. The lost-wake race is
closed by a consume-on-wake flag re-checked after waker registration.

\section{Sitas integration --- the unified shard wait}

A sitas \texttt{ShardExecutor} runs one shard's tasks. It polls ready tasks
within a budget, then blocks in a single \texttt{ReactorBackend::wait} ---
which maps to \texttt{CQ\_WAIT} on the shard's queue. Three event sources
converge on that one wait point:

\begin{itemize}
    \item Kernel completions (CQ entries)
    \item Endpoint readiness (coalesced wake from a CQ-bound endpoint)
    \item Device interrupts (coalesced wake from a bound interrupt object)
\end{itemize}

When the wait returns, the executor drains CQ entries, wakes tasks whose wakers
were registered for the returned handles, and re-enters the polling loop. Task
wakers mark the task ready \emph{and} wake the reactor, so a cross-shard message
injected via a waker-integrated channel releases the target shard's wait and
re-polls exactly the affected task.

\section{Userspace driver model}

A driver protection domain receives a \texttt{DriverGrant} at spawn time:
exactly the MMIO register window and interrupt it needs, delivered as
capabilities through the config-page contract. The driver maps the MMIO region
into its own address space as EL0 device memory, binds the interrupt to its
completion queue, and serves protocol requests from a single \texttt{CQ\_WAIT}.

The reference UART driver demonstrates the full lifecycle:
\begin{itemize}
    \item Direct EL0 MMIO writes to the PL011 transmit FIFO.
    \item Interrupt-driven deferred reads: the driver retains a reply token
        until a device interrupt arrives, then reads the receive register and
        completes the reply.
    \item Cooperative shutdown: unmaps the MMIO region and closes the interrupt.
    \item Crash recovery: an uncooperative exit (the driver terminates without
        releasing its device caps) is followed by teardown that reclaims the
        MMIO and interrupt route, reconciles outstanding operations, and
        invalidates stale connections. A restarted instance with fresh grants
        resumes under a bumped generation.
\end{itemize}

\section{Networking}

The networking architecture treats Ethernet as a generic frame transport. The
NIC driver exports raw frames through \texttt{OP\_SEND} (transmit via a moved
memory object) and \texttt{OP\_RECV} (deferred receive); it knows nothing about
IP or TCP. Higher-level services (reliable message layer, RPC, TCP/IP
compatibility) consume frames through the driver's endpoint.

A smoltcp 0.13 adapter implements the \texttt{phy::Device} trait over the NIC
driver endpoint: \texttt{receive()} maps to \texttt{OP\_RECV},
\texttt{transmit()} maps to \texttt{OP\_SEND} with a memory object. The TCP/IP
stack itself needs no modification. A TCP/IP service binary registers as
\texttt{"tcpip"} and enters a cooperative poll loop.

Protocol crates (\texttt{charlotte-protocol-net},
\texttt{charlotte-protocol-msg}) define the wire formats for frame transport and
reliable sequenced/acked message delivery.

\section{The logical processor}

A \textbf{Logical Processor} (LP) is one hardware core. The kernel shards its
state per-LP using \texttt{PerLp<T>} (lock-wrapped) or
\texttt{ShardLocal<T>} (lock-free, closure-based). LP placement is not advisory:
an LP \emph{is} a core.

\section{The scheduler}

The system scheduler (\texttt{SYSTEM\_SCHEDULER}) uses pluggable
\texttt{LpScheduler} instances. The \texttt{RoundRobin} implementation queues
threads in a FIFO with a per-LP quantum timer. Key operations:
\texttt{spawn\_thread}, \texttt{block\_thread}, \texttt{submit\_ready\_thread},
\texttt{yield\_lp}, \texttt{abort}.

Wakes are idempotent: re-admitting an already-Ready/Running thread is a benign
no-op (the unified shard wait aggregates multiple wake sources onto one thread).

\endinput
